An implementation of this kind must be evaluated not only by its
internal efficiency but also by how coherently it exposes its
capabilities to the user. This is especially important on MCU-based
platforms, where debugging is often more difficult than on Linux-class
systems because the developer has less runtime introspection, fewer
interactive tools, and a much narrower margin for trial-and-error when
something goes wrong during boot or during a real-time task interaction.
For this reason, SentAI does not expose a single debugging path, but a
small set of complementary connection modes over USB that make the board
significantly easier to inspect and operate in practice.

The first mode is an interactive MicroPython REPL available over the
serial interface, which allows a researcher or operator to connect
directly to the board, inspect the runtime state, execute commands live,
load scripts, and request contextual help from within the system itself.
This is a consequential design choice: the platform is not only flashed
and executed as a fixed firmware image, but can also be interrogated and
steered interactively during development and testing. In practice, the
REPL becomes the primary entry point for experimentation, rapid
diagnosis, and iterative mission design.

The second mode is USB networking over IP, through which the firmware
exposes an onboard web server and file-browser interface. This gives the
board a lightweight service surface that can be reached directly from
the host, making it possible to inspect and transfer artifacts without
depending exclusively on the serial console. The third mode is USB mass
storage, in which the LittleFS user partition is exposed as a mounted
drive so that images, scripts, and models can be uploaded or downloaded
directly from the host machine. Taken together, these three USB-facing
paths, namely serial REPL, IP connectivity with web access, and mounted
LittleFS storage, form a practical answer to the usual debugging and
deployment friction of MCU-based systems.

The built-in help interface available from the REPL makes this
programming model discoverable without forcing the user to leave the
target system. This matters because the runtime is relatively rich: it
does not expose a single detector call or a single flight command, but a
full namespace of platform primitives intended to support sensing,
inference, telemetry, storage, and high-level control from one embedded
environment. The same environment also supports script-driven startup
through \texttt{/main.py}, which is automatically executed at boot
before the interactive session begins. This gives the platform an
important dual character: it can behave as an interactive research
instrument during development, but it can also behave as a self-starting
embedded application once the desired mission logic has been written.

This startup behavior is complemented by a boot logging mechanism
implemented directly in the runtime. During initialization, console and
driver output are captured into \texttt{/log/boot.log}, while the
previous boot log is rotated to \texttt{/log/boot\_old.log}. This is an
important practical feature for embedded debugging, because many
failures on an MCU platform occur before a user can attach interactively
to the REPL. By preserving the initialization trace on LittleFS, the
system allows the user to inspect the outcome of FreeRTOS bring-up,
early peripheral initialization, and startup script execution even after
the board has already moved past the failing point.

The \texttt{sentai} namespace organizes the platform into functional
modules such as \texttt{io}, \texttt{rtos}, \texttt{tpu}, \texttt{fs},
\texttt{camera}, \texttt{imu}, \texttt{usb}, \texttt{link},
\texttt{crazy}, and \texttt{pipeline}, alongside several other utility
interfaces. This structure is significant because it reveals the
intended user abstraction: the operator is not expected to manipulate a
collection of unrelated firmware endpoints, but to work with a coherent
runtime in which device functions are grouped by task and by
experimental purpose. The platform can query RTOS state, inspect heap
usage, load models, capture or switch cameras, read inertial
measurements, start a continuous detector, track objects, communicate
over MAVLink, or command a drone, all through a single embedded language
environment.

This helps explain why MicroPython is a central implementation choice
rather than a peripheral convenience. The interpreter does not replace
the lower-level real-time system; instead, it defines a stable
experimentation layer above it. The help text demonstrates this duality
clearly. On one hand, users access simple operations such as
\texttt{sentai.camera.to\_tensor()} or \texttt{sentai.tpu.invoke()}. On
the other hand, the same interface reaches into advanced platform
features such as tracking events, per-camera geometry, ground-plane
projection, and MAVLink message exchange. The implementation thereby
condenses a complex firmware system into a form that is operable from
concise Python scripts.

Several modules are especially important for understanding the research
value of the platform. The \texttt{rtos} module exposes runtime
observability primitives such as task inspection, CPU statistics, and
heap information, which are essential when the system is tuned for
sustained real-time throughput. The \texttt{tpu} module provides the
basic inference control plane: model loading, tensor access, invocation,
quantization metadata, and detector-oriented post-processing. The
\texttt{camera} module complements it by managing capture, JPEG export,
transfer to the tensor path, runtime camera switching, and frame
sequencing. The \texttt{imu} module supplies the inertial measurements
needed by viewpoint-aware tracking and projection. The \texttt{fs} and
\texttt{usb} modules support a practical workflow in which models,
scripts, captured outputs, and debugging artifacts can be transferred to
and from the device without rebuilding the firmware for every iteration.

The help surface also shows that these capabilities were designed to be
exercised directly from the target rather than only through host-side
tooling. Typical examples include immediate tensor feeding from the
camera, live detector startup, autopilot connection, or direct Crazyflie
flight commands, for example:

\begin{Shaded}
\begin{Highlighting}[]
\NormalTok{sentai.camera.to\_tensor()}
\NormalTok{sentai.tpu.invoke()}
\BuiltInTok{print}\NormalTok{(sentai.tpu.detect())}
\end{Highlighting}
\end{Shaded}

\begin{Shaded}
\begin{Highlighting}[]
\NormalTok{sentai.link.init(}\DecValTok{57600}\NormalTok{)}
\NormalTok{sentai.pipeline.start(}\FloatTok{0.3}\NormalTok{, }\FloatTok{0.45}\NormalTok{, }\DecValTok{50}\NormalTok{, }\VariableTok{True}\NormalTok{)}
\ControlFlowTok{while} \VariableTok{True}\NormalTok{:}
\NormalTok{    msg }\OperatorTok{=}\NormalTok{ sentai.link.receive()}
    \ControlFlowTok{if}\NormalTok{ msg:}
        \BuiltInTok{print}\NormalTok{(msg)}
\end{Highlighting}
\end{Shaded}

\begin{Shaded}
\begin{Highlighting}[]
\NormalTok{sentai.crazy.init()}
\NormalTok{sentai.crazy.arm()}
\NormalTok{sentai.crazy.fly(}\FloatTok{0.3}\NormalTok{, }\DecValTok{3000}\NormalTok{)}
\NormalTok{sentai.crazy.fly\_stop()}
\end{Highlighting}
\end{Shaded}

At the control and autonomy layer, \texttt{link}, \texttt{crazy}, and
\texttt{pipeline} form the most distinctive part of the namespace. The
\texttt{link} module exposes MAVLink communication as a first-class
runtime capability, allowing the board to exchange telemetry and
structured messages with an autopilot or a ground-control system. The
\texttt{crazy} module does the same for the indoor Crazyflie target,
exposing high-level flight functions directly to Python. The
\texttt{pipeline} module unifies the continuous detector, tracker, event
stream, per-camera geometry, and pose-aware projection logic into a
single interface. Together, these modules reveal the real ambition of
the implementation: the embedded AI platform is not only an inference
endpoint, but a scriptable autonomy substrate in which perception,
communication, and high-level flight behavior can be composed
interactively.

The \texttt{pipeline} module is especially illustrative. According to
the documented runtime surface, it exposes continuous detection,
statistics, tracker snapshots, event streams, configurable association
parameters, pose setting, and per-camera geometry. This is not a trivial
wrapper around a detector. It is the public face of the deeper
implementation presented in the previous subsections: the parallel
camera-to-TPU pipeline, the tracking-by-detection logic, and the
geometry-aware projection model all become accessible through a minimal
but expressive Python API. The same can be said for the \texttt{link}
and \texttt{crazy} modules, which reveal that flight and telemetry
functions are not separate applications but first-class capabilities of
the runtime itself.

From a research-methodology perspective, this matters because it changes
the pace and style of experimentation. A researcher can move from
firmware bring-up to mission logic, data capture, tracking inspection,
and communication testing without leaving the same runtime environment.
This reduces iteration cost and makes the platform more suitable for
exploratory autonomy research, where control policies, perception
thresholds, and communication strategies change frequently.

The built-in help surface also provides indirect evidence of
implementation maturity. A platform that documents precise operational
behaviors, module boundaries, argument conventions, and example
workflows has generally progressed beyond proof-of-concept code. In
SentAI, the fact that these capabilities can be explored directly from
the serial REPL supports the claim that the implementation has been
organized as a reusable experimental environment rather than as a
collection of hidden firmware entry points.

The programming model therefore completes the overall argument of the
chapter. The EdgeTPU provides the inference power, the MCU and RTOS
provide the orchestration, the cameras and IMU provide the sensing
context, and MicroPython makes the whole system operable as a research
platform. The implementation succeeds not only when it runs efficiently,
but when it turns that efficiency into a usable and extensible interface
for autonomous drone experimentation.
